% FILE: chapters/ch06_s4_llms.tex
\section{Large Language Models}\label{sec:llms}

In Section~\ref{sec:nlp}, we studied foundational Natural Language Processing techniques: tokenization, dense semantic embeddings, and training task-specific text classifiers. While those models excel at dedicated narrow tasks (such as categorizing an industrial maintenance report into ``normal'' vs ``critical''), they lack general linguistic understanding, logical reasoning, and zero-shot problem-solving capabilities.

Over the past several years, a transformative paradigm shift has reshaped the artificial intelligence landscape: the emergence of \textbf{Large Language Models (LLMs)}. Pioneered by models such as GPT-4, LLaMA, Gemini, and Mistral, LLMs scale the decoder-only Transformer architecture (introduced in Section~\ref{sec:attention_transformers}) to billions or trillions of parameters, trained across massive multi-terabyte web-scale corpora.

In modern Internet of Things (IoT) ecosystems, industrial robotics, and cyber-physical engineering, LLMs and their compact counterparts---\textbf{Small Language Models (SLMs)}---are rapidly becoming the central cognitive orchestrators:
\begin{itemize}
    \item \textbf{Autonomous Edge Agents}: An LLM ingests multi-sensor telemetry, correlates unexpected thermal spikes with historical maintenance logs, diagnoses an emerging mechanical seizure, and autonomously issues corrective commands across an MQTT message bus.
    \item \textbf{Natural Language Interfaces for Actuators}: Human plant operators issue high-level instructions (e.g., \emph{``Reroute excess solar generation to charge battery bank two and reduce chiller power by twenty percent''}), which the language model translates into structured, validated JSON payloads that control physical relays and programmable logic controllers (PLCs).
    \item \textbf{Diagnostic Retrieval over Technical Manuals}: Correlating obscure error codes (e.g., \code{"ERROR_CAN_0x4F"}) with thousands of pages of PDF equipment schematics in seconds via Retrieval-Augmented Generation (RAG).
\end{itemize}

In this section, we study how LLMs are trained, master prompt engineering and structured generation, demystify Retrieval-Augmented Generation (RAG), examine parameter-efficient fine-tuning with LoRA, explore objective evaluation metrics, and learn how to run quantized models locally on edge gateways.

% ==============================================================================
% 1. HOW LLMS ARE TRAINED
% ==============================================================================
\subsection{How LLMs Are Trained: The Three-Stage Alignment Pipeline}\label{subsec:llm_training}

Building a modern, instruction-following language model is not a single training run. It is a rigorous three-stage pipeline:

\begin{figure}[htbp]
\centering
\begin{tikzpicture}[
    scale=0.88,
    stage/.style={rectangle, draw=keyideablue, fill=keyideablue!15, rounded corners=2mm, thick, minimum width=3.4cm, minimum height=1.6cm, align=center},
    arrow/.style={-{Stealth[length=2.5mm]}, thick, draw=gray!80}
]
    % Stage 1
    \node[stage] (s1) at (0, 0) {
        \textbf{Stage 1: Pretraining}\\
        \scriptsize Trillions of raw tokens\\
        \scriptsize Next-token prediction\\
        \scriptsize Creates Base Model
    };

    % Stage 2
    \node[stage, draw=analogygreen, fill=analogygreen!15] (s2) at (4.6, 0) {
        \textbf{Stage 2: SFT}\\
        \scriptsize Instruction Tuning\\
        \scriptsize Curated Q\&A pairs\\
        \scriptsize Teaches dialog format
    };

    % Stage 3
    \node[stage, draw=warningred, fill=warningred!15] (s3) at (9.2, 0) {
        \textbf{Stage 3: Alignment}\\
        \scriptsize RLHF / DPO\\
        \scriptsize Human preferences\\
        \scriptsize Helpful, Honest, Harmless
    };

    % Connecting arrows
    \draw[arrow] (s1) -- (s2);
    \draw[arrow] (s2) -- (s3);
\end{tikzpicture}
\caption{The three-stage training pipeline: Self-supervised Pretraining on raw text, Supervised Fine-Tuning (SFT) on instructions, and Preference Alignment (RLHF/DPO).}\label{fig:llm_pipeline}
\end{figure}

\subsubsection{Stage 1: Self-Supervised Pretraining (The Foundation)}
The model starts with randomly initialized weights. It is fed trillions of tokens of raw text scraped from books, code repositories, research papers, and technical manuals.
\begin{itemize}
    \item \textbf{Training Objective}: \textbf{Autoregressive Next-Token Prediction}. Given a sequence of tokens $x_1, x_2, \dots, x_{t-1}$, the network predicts the probability distribution of the immediate next token $x_t$:
    \begin{equation}
    \mathcal{L}_{\text{pretrain}}(\theta) = -\sum_{t=1}^T \log P(x_t \mid x_1, x_2, \dots, x_{t-1}; \theta)\label{eq:pretrain_loss}
    \end{equation}
    \item \textbf{The Result}: A \textbf{Base Model}. The base model has acquired broad linguistic grammar, facts about the world, and code syntax. However, it is \emph{not} a helpful assistant. If you ask it: \code{"What causes pump cavitation?"}, it might simply continue your text with: \code{"What causes motor overheating? Question 3..."} because it treats your question as an exam paper to be completed, not an instruction to be answered!
\end{itemize}

\subsubsection{Stage 2: Supervised Fine-Tuning (Instruction Tuning)}
To transform a raw completion engine into an interactive assistant, we fine-tune the base model on hundreds of thousands of meticulously curated \textbf{Instruction-Response Pairs}:
\begin{equation*}
\left( \text{User: ``Explain what caused the temperature spike.''}, \text{Assistant: ``The spike was caused by...''} \right)
\end{equation*}
The model learns the conversational roles (\texttt{<|user|>}, \texttt{<|assistant|>}) and learns to directly answer questions, summarize documents, and generate code.

\subsubsection{Stage 3: Preference Alignment (RLHF and DPO)}
Even after instruction tuning, models can hallucinate, produce toxic or dangerous advice, or ramble incoherently. \textbf{Alignment} aligns model outputs with the ``Three Hs'': \textbf{Helpful, Honest, and Harmless}.
\begin{itemize}
    \item \textbf{RLHF (Reinforcement Learning from Human Feedback)}: Human evaluators score multiple model candidate responses. A separate \textbf{Reward Model} learns human preferences, and the language model is updated via Reinforcement Learning (PPO: Proximal Policy Optimization) to maximize the reward.
    \item \textbf{DPO (Direct Preference Optimization)}: A modern, simpler formulation (Rafailov et al., 2023) that directly optimizes the language model on pairs of preferred ($y_w$) and dispreferred ($y_l$) responses without training a separate reward model or using complex reinforcement learning:
    \begin{equation}
    \mathcal{L}_{\text{DPO}}(\theta) = -\mathbb{E}_{(x, y_w, y_l)}\left[ \log \sigma\left( \beta \log \frac{\pi_\theta(y_w \mid x)}{\pi_{\text{ref}}(y_w \mid x)} - \beta \log \frac{\pi_\theta(y_l \mid x)}{\pi_{\text{ref}}(y_l \mid x)} \right) \right]\label{eq:dpo_loss}
    \end{equation}
    where $\pi_{\text{ref}}$ is the frozen SFT reference model and $\beta$ controls drift from the baseline.
\end{itemize}

% ==============================================================================
% 2. PROMPT ENGINEERING
% ==============================================================================
\subsection{Prompt Engineering: Steering Generative Engines}\label{subsec:prompt_engineering}

Because LLMs are probabilistic completion engines, how an input is formatted profoundly alters the model's sampling trajectory. \textbf{Prompt Engineering} is the practice of designing, structuring, and optimizing inputs to reliably guide an LLM toward accurate, predictable outputs.

\paragraph{1. Zero-Shot vs Few-Shot Prompting (In-Context Learning)}
\begin{itemize}
    \item \textbf{Zero-Shot}: You provide only the instruction without any examples:
    \begin{lstlisting}[style=outputstyle]
Classify the sensor anomaly: "Vibration frequency surged to 120 Hz at 80C."
    \end{lstlisting}
    \item \textbf{Few-Shot Prompting}: You provide 2 to 5 input-output demonstration pairs within the prompt before providing the target query. The model leverages its \textbf{in-context learning} capacity to mimic the pattern:
    \begin{lstlisting}[style=outputstyle]
Log: "Current jumped to 45A instantly." -> Label: ELECTRICAL_SURGE
Log: "Bearing temperature rose 1C every hour." -> Label: THERMAL_DRIFT
Log: "Vibration frequency surged to 120 Hz at 80C." -> Label:
    \end{lstlisting}
\end{itemize}

\paragraph{2. Chain-of-Thought (CoT) Prompting}
Introduced by Jason Wei et al. (2022), \textbf{Chain-of-Thought (CoT)} prompting encourages the model to generate intermediate reasoning steps before arriving at a final answer:
\begin{itemize}
    \item By appending the simple trigger phrase: \code{"Let's think step by step"}, the model allocates extra output tokens to compute logical deductions.
    \item This drastically reduces arithmetic errors and hallucinated diagnostic conclusions in industrial IoT troubleshooting.
\end{itemize}

\paragraph{3. Structured Output Enforcement (JSON for IoT)}
In embedded engineering, an LLM's output must frequently interface with downstream code (such as publishing an MQTT topic or controlling an ESP32 relay). A conversational response like \emph{``Sure, I'd be happy to turn off the valve for you!''} breaks code parsers.
\begin{itemize}
    \item Modern inference frameworks (such as llama.cpp and OpenAI JSON mode) support \textbf{Grammar-Constrained Sampling}: the model's output logits are masked at each token step to enforce strict conformance to a predefined JSON schema:
    \begin{lstlisting}[style=outputstyle]
{
  "target_device": "solenoid_valve_3",
  "action": "CLOSE",
  "override_safety": false
}
    \end{lstlisting}
\end{itemize}

% ==============================================================================
% 3. RETRIEVAL-AUGMENTED GENERATION (RAG)
% ==============================================================================
\subsection{Retrieval-Augmented Generation (RAG)}\label{subsec:rag}

\begin{keyidea}[RAG: The Open-Book Exam]
An LLM's parametric memory is frozen at the moment training ends. It suffers from \textbf{knowledge cutoff} (unaware of recent events) and \textbf{hallucination} (confidently inventing plausible-sounding but false technical specifications).

\textbf{Retrieval-Augmented Generation (RAG)} transforms the LLM into an \emph{open-book test taker}:
\begin{enumerate}
    \item When a query arrives, a search engine retrieves authoritative text excerpts from a private knowledge base (e.g., equipment manuals, maintenance tickets).
    \item The retrieved context is injected directly into the LLM's prompt.
    \item The LLM synthesizes its answer based \emph{strictly on the retrieved facts}, citing exact sources.
\end{enumerate}
\end{keyidea}

\begin{figure}[htbp]
\centering
\begin{tikzpicture}[
    scale=0.85,
    block/.style={rectangle, draw=keyideablue, fill=keyideablue!15, rounded corners=1.5mm, thick, minimum width=2.6cm, minimum height=0.9cm, font=\scriptsize\sffamily\bfseries, align=center},
    db/.style={cylinder, draw=analogygreen, fill=analogygreen!20, shape border rotate=90, aspect=0.25, thick, minimum width=1.8cm, minimum height=1.4cm, font=\scriptsize\sffamily\bfseries, align=center},
    io/.style={rectangle, draw=gray!70, fill=white, rounded corners=1mm, thick, minimum width=2.0cm, minimum height=0.6cm, font=\scriptsize\sffamily},
    arrow/.style={-{Stealth[length=2.5mm]}, thick, draw=gray!80}
]
    % Left: Ingestion Phase
    \node[io] (docs) at (0, 3.2) {PDF Manuals / Logs};
    \node[block] (chunk) at (0, 1.8) {Chunking \&\\Embedding Model};
    \node[db] (vdb) at (0, -0.6) {Vector DB\\(Chroma/FAISS)};
    \draw[arrow] (docs) -- (chunk);
    \draw[arrow] (chunk) -- node[left, font=\tiny] {Dense Vectors} (vdb);

    % Right: Query & Generation Phase
    \node[io] (query) at (5.0, 3.2) {User / Sensor Query};
    \node[block] (q_emb) at (5.0, 1.8) {Embed Query};
    \node[block, fill=warningred!15, draw=warningred] (prompt_build) at (5.0, -0.6) {Augmented Prompt\\$\text{Context} + \text{Query}$};
    \node[block, fill=keyideablue!25] (llm) at (8.8, -0.6) {LLM Generator};
    \node[io, fill=analogygreen!20] (answer) at (8.8, 1.8) {Grounded Answer\\+ Citations};

    % Connections
    \draw[arrow] (query) -- (q_emb);
    \draw[arrow] (q_emb) -- (vdb);
    \draw[arrow, analogygreen!80!black, line width=1.2pt] (vdb) -- node[above, font=\tiny\bfseries] {Top-$K$ Chunks} (prompt_build);
    \draw[arrow] (q_emb) -- (prompt_build);
    \draw[arrow] (prompt_build) -- (llm);
    \draw[arrow] (llm) -- (answer);
\end{tikzpicture}
\caption{The Retrieval-Augmented Generation (RAG) Architecture. Private manuals are embedded and stored in a vector database; at query time, relevant text chunks are dynamically retrieved and supplied to the LLM.}\label{fig:rag_pipeline}
\end{figure}

\paragraph{The Three Steps of RAG}
\begin{enumerate}
    \item \textbf{Ingestion and Indexing}:
    \begin{itemize}
        \item Long technical documents are split into smaller \textbf{chunks} (typically 256 to 512 tokens with 50-token overlap).
        \item An embedding model (such as \code{text-embedding-3} or \code{bge-small}) converts each chunk into a dense vector $\mathbf{v} \in \mathbb{R}^d$.
        \item Vectors and text chunks are indexed in a \textbf{Vector Database} (e.g., FAISS, Chroma, Qdrant).
    \end{itemize}
    \item \textbf{Retrieval}:
    \begin{itemize}
        \item When the technician submits a query $\mathbf{q}$ (e.g., \emph{``What is the maximum torque for bolt M12 on motor bracket?''}), the query is converted into an embedding $\mathbf{v}_q$.
        \item The vector database computes the cosine similarity (Equation~\eqref{eq:cosine_similarity}) between $\mathbf{v}_q$ and all indexed chunks, retrieving the \textbf{Top-$K$} most similar passages.
    \end{itemize}
    \item \textbf{Augmented Generation}:
    \begin{itemize}
        \item The retrieved excerpts are inserted into an augmented prompt:
        \begin{lstlisting}[style=outputstyle]
Context: [Manual Excerpt: Section 4.2 states M12 bracket bolts must be torqued to 85 Nm.]
Question: What is the maximum torque for bolt M12 on motor bracket?
Answer using ONLY the provided context:
        \end{lstlisting}
        \item The LLM generates: \code{"The maximum torque for bolt M12 is 85 Nm (Section 4.2)."}.
    \end{itemize}
\end{enumerate}

% ==============================================================================
% 4. PARAMETER-EFFICIENT FINE-TUNING (LORA)
% ==============================================================================
\subsection{Parameter-Efficient Fine-Tuning: LoRA}\label{subsec:lora}

\paragraph{The VRAM Barrier of Full Fine-Tuning}
A modern 8-billion parameter model (such as LLaMA-3-8B) requires approximately 16 GB of memory just to store its weights in 16-bit precision. During full backpropagation, storing optimizer states (Adam stores two 32-bit floats per parameter), gradients, and activations requires \textbf{over 80 GB of GPU VRAM}, requiring multi-thousand-dollar server clusters. Furthermore, full fine-tuning generates a complete 16 GB weight checkpoint for every customized task.

\begin{keyidea}[Low-Rank Adaptation (LoRA)]
In 2021, Edward Hu and colleagues at Microsoft introduced \textbf{LoRA (Low-Rank Adaptation)}. 

LoRA is based on a profound empirical insight: \textbf{the weight updates $\Delta \mathbf{W}$ that occur during task adaptation have a very low ``intrinsic dimension'' (rank)}. 

Instead of updating the full $d \times k$ weight matrix, LoRA freezes the original pretrained weights $\mathbf{W}_0$ and decomposes the update into the product of two tiny, low-rank matrices: $\Delta \mathbf{W} = \mathbf{B} \mathbf{A}$, where rank $r \ll \min(d, k)$.
\end{keyidea}

\paragraph{Mathematical Formulation of LoRA}
For a frozen linear layer $\mathbf{W}_0 \in \mathbb{R}^{d \times k}$ (such as an attention projection matrix $\mathbf{W}^Q$ or $\mathbf{W}^V$):
\begin{equation}
\mathbf{h} = \mathbf{W}_0 \mathbf{x} + \Delta \mathbf{W} \mathbf{x} = \mathbf{W}_0 \mathbf{x} + \frac{\alpha}{r} (\mathbf{B} \mathbf{A}) \mathbf{x}\label{eq:lora_eq}
\end{equation}
where:
\begin{itemize}
    \item $\mathbf{W}_0 \in \mathbb{R}^{d \times k}$ is the \textbf{frozen} pretrained weight matrix. No gradients are computed for it.
    \item $\mathbf{A} \in \mathbb{R}^{r \times k}$ is a trainable adapter matrix initialized with a random Gaussian distribution $\mathcal{N}(0, \sigma^2)$.
    \item $\mathbf{B} \in \mathbb{R}^{d \times r}$ is a trainable adapter matrix initialized to \textbf{all zeros}.
    \item Because $\mathbf{B} = \mathbf{0}$ at initialization, $\Delta \mathbf{W} = \mathbf{B}\mathbf{A} = \mathbf{0}$. The model starts with its exact pretrained behavior!
    \item $\frac{\alpha}{r}$ is a constant scaling factor (where $\alpha$ is a hyperparameter, typically set to $2r$).
\end{itemize}

\begin{figure}[htbp]
\centering
\begin{tikzpicture}[
    scale=0.9,
    block/.style={rectangle, draw=gray!70, fill=gray!15, thick, minimum width=2.4cm, minimum height=2.8cm, font=\small\sffamily\bfseries, align=center},
    lora_a/.style={rectangle, draw=warningred, fill=warningred!20, thick, minimum width=2.4cm, minimum height=0.6cm, font=\scriptsize\sffamily\bfseries},
    lora_b/.style={rectangle, draw=warningred, fill=warningred!20, thick, minimum width=0.8cm, minimum height=2.0cm, font=\scriptsize\sffamily\bfseries},
    op/.style={circle, draw=keyideablue, fill=keyideablue!20, thick, minimum size=0.7cm, font=\large\bfseries},
    arrow/.style={-{Stealth[length=2.5mm]}, thick, draw=gray!80}
]
    % Input
    \node (x) at (0, -2.4) {Input $\mathbf{x} \in \mathbb{R}^k$};

    % Frozen Pretrained Weights (Left path)
    \node[block] (w0) at (-2.0, 0) {Frozen Pretrained\\$\mathbf{W}_0 \in \mathbb{R}^{d \times k}$\\\scriptsize(No Gradients)};

    % Trainable LoRA Matrices (Right path)
    \node[lora_a] (a) at (2.2, -0.8) {Matrix $\mathbf{A} \in \mathbb{R}^{r \times k}$};
    \node[lora_b] (b) at (2.2, 0.8) {$\mathbf{B} \in \mathbb{R}^{d \times r}$};

    % Addition operator
    \node[op] (add) at (0, 2.4) {$+$};
    \node (out) at (0, 3.6) {Output $\mathbf{h} = \mathbf{W}_0\mathbf{x} + \frac{\alpha}{r}\mathbf{B}\mathbf{A}\mathbf{x}$};

    % Connections
    \draw[arrow] (x) -- (0, -1.8) -| (w0.south);
    \draw[arrow] (x) -- (0, -1.8) -| (a.south);
    \draw[arrow] (a) -- node[right, font=\tiny] {$r$} (b);
    \draw[arrow] (w0.north) |- (add.west);
    \draw[arrow] (b.north) |- (add.east);
    \draw[arrow] (add) -- (out);
\end{tikzpicture}
\caption{The LoRA Architecture. Pretrained weight matrix $\mathbf{W}_0$ is frozen. Trainable rank decomposition matrices $\mathbf{A}$ and $\mathbf{B}$ capture task-specific adaptations with over $99\%$ fewer parameters.}\label{fig:lora_arch}
\end{figure}

\begin{example}[Calculating LoRA Parameter Savings]
Consider a standard Transformer projection matrix with dimension $d = 4{,}096$ and $k = 4{,}096$. We attach a LoRA adapter with rank $r = 8$.

\textbf{1. Full Fine-Tuning Parameter Count}:
\begin{equation*}
\text{Params}_{\text{full}} = 4{,}096 \times 4{,}096 = \mathbf{16{,}777{,}216 \text{ parameters}}
\end{equation*}

\textbf{2. LoRA Parameter Count}:
\begin{align*}
\text{Matrix } \mathbf{A} \in \mathbb{R}^{8 \times 4096} &\implies 8 \times 4{,}096 = 32{,}768 \text{ parameters} \\
\text{Matrix } \mathbf{B} \in \mathbb{R}^{4096 \times 8} &\implies 4{,}096 \times 8 = 32{,}768 \text{ parameters} \\
\text{Total LoRA Params} &= 32{,}768 + 32{,}768 = \mathbf{65{,}536 \text{ parameters}}
\end{align*}

\textbf{3. Parameter Reduction}:
\begin{equation*}
\text{Reduction Ratio} = \frac{65{,}536}{16{,}777{,}216} = \frac{1}{256} \approx \mathbf{0.39\%}
\end{equation*}
LoRA tunes only \textbf{$0.39\%$ of the original parameters} (a $99.61\%$ reduction!). Instead of a 16 GB checkpoint, the saved LoRA adapter file weighs only \textbf{131 KB}! Multiple specialized adapters (e.g., one for pump troubleshooting, one for electrical wiring) can be swapped in RAM dynamically in milliseconds.
\end{example}

% ==============================================================================
% 5. EVALUATING LLM OUTPUTS
% ==============================================================================
\subsection{Evaluating LLM Outputs}\label{subsec:llm_evaluation}

Evaluating generative language models is fundamentally more challenging than evaluating classification models. In classification, a prediction is either right or wrong. In text generation, there are hundreds of valid ways to phrase a diagnosis.

\paragraph{1. Perplexity (PPL)}
Perplexity measures how well a probability model predicts a sample of text. It is the exponentiated cross-entropy loss:
\begin{equation}
\operatorname{PPL} = \exp\left( \mathcal{L} \right) = \exp\left( -\frac{1}{T} \sum_{t=1}^T \ln P(x_t \mid x_{<t}) \right)\label{eq:perplexity}
\end{equation}
A lower perplexity indicates the model is less surprised by the text, reflecting stronger linguistic modeling capability.

\paragraph{2. N-Gram Overlap Metrics: BLEU and ROUGE}
\begin{itemize}
    \item \textbf{ROUGE-L}: Measures the Longest Common Subsequence (LCS) between the generated text and a reference answer. It assesses recall (did the model capture the essential reference facts?).
    \item \emph{Limitation}: ROUGE relies on exact string matches. If the reference says \code{"shut down the pump"} and the model generates \code{"deactivate the pump"}, ROUGE penalizes the model despite the two phrases being semantically identical!
\end{itemize}

\paragraph{3. LLM-as-a-Judge and RAG Triad Metrics}
In modern practice, engineers evaluate LLMs using stronger foundation models (such as GPT-4) guided by strict rubric criteria:
\begin{enumerate}
    \item \textbf{Faithfulness}: Did the model invent claims not supported by the retrieved context? (Detects hallucinations).
    \item \textbf{Answer Relevance}: Did the generated response directly address the user's specific question?
    \item \textbf{Context Precision}: Were the retrieved manual excerpts actually relevant to the query?
\end{enumerate}

% ==============================================================================
% 6. CODE: BUILDING A VECTOR RAG PIPELINE
% ==============================================================================
\subsection{Python Implementation: A Self-Contained Vector RAG Pipeline}\label{subsec:python_rag}

Let us build a complete, runnable vector retrieval pipeline in Python and PyTorch. We will index equipment manuals and retrieve the exact diagnostic procedure for an incoming machine error code.

\paragraph{The Python Implementation}
\begin{lstlisting}[language=Python]
import torch
import torch.nn.functional as F

# 1. Technical Knowledge Base (Equipment Maintenance Chunks)
knowledge_base = [
    "Doc 0: Motor bearing wear is indicated by high-frequency vibration above 100 Hz. Action: Replace bearing.",
    "Doc 1: Pump cavitation occurs when suction pressure drops below 1.2 bar. Action: Open suction valve.",
    "Doc 2: Hydraulic fluid over-temperature exceeds 80C. Action: Check cooling fan and fluid radiator.",
    "Doc 3: Inverter CAN-bus error 0x4F indicates disconnected termination resistor. Action: Inspect wiring."
]

# 2. Simulated Embedding Model (Bag-of-Words Projection for Reproducibility)
# Maps documents to normalized semantic vectors
toy_vocab = {
    "vibration": 0, "bearing": 1, "cavitation": 2, "pressure": 3,
    "temperature": 4, "hydraulic": 5, "can-bus": 6, "error": 7, "wiring": 8
}

def embed_text(text):
    vec = torch.zeros(len(toy_vocab))
    for word in text.lower().split():
        clean_word = word.strip(".:,")
        if clean_word in toy_vocab:
            vec[toy_vocab[clean_word]] += 1.0
    # Normalize to unit vector for cosine similarity
    norm = torch.norm(vec, p=2)
    return vec / norm if norm > 0 else vec

# Index all documents in the knowledge base
doc_embeddings = torch.stack([embed_text(doc) for doc in knowledge_base])

# 3. User Query (Arriving from an industrial operator)
query = "We observed severe cavitation and low suction pressure on pump 2."
query_embedding = embed_text(query)

# 4. Vectorized Cosine Similarity Search: Q @ Docs^T
similarity_scores = torch.matmul(doc_embeddings, query_embedding)

# Retrieve Top-1 Most Relevant Document
top1_idx = torch.argmax(similarity_scores).item()
top1_score = similarity_scores[top1_idx].item()

print(f"User Query:      {query}")
print(f"Top Match Doc:   #{top1_idx} (Similarity: {top1_score:.4f})")
print(f"Retrieved Text:  {knowledge_base[top1_idx]}")

# 5. Synthesize Augmented Prompt for LLM Generation
augmented_prompt = f"""[SYSTEM: Answer the query using ONLY the provided context below.]
CONTEXT:
{knowledge_base[top1_idx]}

QUERY:
{query}

ACTIONABLE DIAGNOSIS:"""

print("\n--- Synthesized RAG Prompt ---")
print(augmented_prompt)
\end{lstlisting}

\paragraph{Expected Output}
\begin{lstlisting}[style=outputstyle]
User Query:      We observed severe cavitation and low suction pressure on pump 2.
Top Match Doc:   #1 (Similarity: 0.8165)
Retrieved Text:  Doc 1: Pump cavitation occurs when suction pressure drops below 1.2 bar. Action: Open suction valve.

--- Synthesized RAG Prompt ---
[SYSTEM: Answer the query using ONLY the provided context below.]
CONTEXT:
Doc 1: Pump cavitation occurs when suction pressure drops below 1.2 bar. Action: Open suction valve.

QUERY:
We observed severe cavitation and low suction pressure on pump 2.

ACTIONABLE DIAGNOSIS:
\end{lstlisting}

\paragraph{Line-by-Line Code Breakdown}
\begin{itemize}
    \item \textbf{Lines 5--10}: We define a realistic private knowledge base containing technical diagnostic rules for motor and pump maintenance.
    \item \textbf{Line 24}: \code{norm = torch.norm(vec, p=2); return vec / norm}. Normalizing vectors to unit length ($\|\mathbf{v}\|_2 = 1$) simplifies cosine similarity: the dot product of two unit vectors is mathematically identical to their cosine similarity!
    \item \textbf{Line 34}: \code{similarity_scores = torch.matmul(doc_embeddings, query_embedding)}. Performs simultaneous, parallel vector comparison across all indexed documents in a single GPU/CPU matrix operation.
    \item \textbf{Lines 37--41}: The system matches the query to Doc 1 with $0.8165$ similarity, rejecting irrelevant motor, hydraulic, and CAN-bus documents.
    \item \textbf{Lines 44--52}: The system constructs the grounded RAG prompt, ensuring the LLM's response cites the exact action (\code{"Open suction valve"}) without hallucination.
\end{itemize}

% ==============================================================================
% 7. IOT & TINYML CONNECTION
% ==============================================================================
\subsection{ML for IoT and Edge: Small Language Models on Edge Gateways}\label{subsec:iot_llm}

\begin{iotbox}[Small Language Models (SLMs) on the Factory Floor]
Deploying language models directly onto local edge gateways (e.g., Raspberry Pi 5, Jetson Orin Nano, industrial DIN-rail IPCs) provides three non-negotiable industrial advantages:
\begin{enumerate}
    \item \textbf{Air-Gapped Privacy}: Industrial telemetry and proprietary assembly blueprints never leave the factory firewall.
    \item \textbf{Zero Cloud Latency}: Local inference responds in milliseconds without waiting for cloud API round-trips.
    \item \textbf{Uninterrupted Offline Reliability}: If external internet connectivity is severed during a storm, the local edge safety system continues operating autonomously.
\end{enumerate}
\end{iotbox}

\paragraph{Quantization and 4-Bit Execution with GGUF}
While a 70-billion parameter model cannot run locally, state-of-the-art \textbf{Small Language Models} (such as LLaMA-3.2-1B, LLaMA-3.2-3B, and Phi-3-mini) achieve remarkable reasoning capabilities. 

To run on edge gateways with limited RAM, models are quantized using \textbf{4-bit Weight Quantization (GGUF / AWQ)}:
\begin{equation*}
\text{Memory Budget} \approx (\text{Parameter Count in Billions}) \times 0.6 \text{ GB}
\end{equation*}
\begin{itemize}
    \item A 1-billion parameter model in 4-bit requires only \textbf{$\approx 600\text{ MB}$ of RAM}!
    \item A 3-billion parameter model requires only \textbf{$\approx 1.8\text{ GB}$ of RAM}!
    \item Running via optimized C++ engines (like \code{llama.cpp}), a Raspberry Pi 5 can generate \textbf{15 to 25 tokens per second}, fast enough for real-time edge voice interaction and autonomous MQTT control.
\end{itemize}

% ==============================================================================
% 8. COMMON MISTAKES AND MISCONCEPTIONS
% ==============================================================================
\subsection{Common Mistakes and Misconceptions}\label{subsec:mistakes}

\begin{warning}[Five Pitfalls in LLM and RAG Systems]
\begin{enumerate}
    \item \textbf{Treating LLMs as Factual Databases}: Expecting a base LLM to remember exact hardware part numbers from memory. LLMs are generative compression models, not databases; they hallucinate facts with supreme confidence. Always ground domain-specific queries with RAG.
    \item \textbf{The ``Lost in the Middle'' Phenomenon}: Packing 50 retrieved chunks into a massive context window. Research (Liu et al., 2023) proves that language models pay strong attention to the beginning and end of their context window, but frequently ignore or ``forget'' facts placed in the middle. Always rank and filter retrieved chunks to provide only the Top-3 to Top-5 most relevant passages.
    \item \textbf{Confusing Fine-Tuning with Knowledge Injection}: Attempting to teach an LLM fresh factual information (such as today's factory inventory) via fine-tuning. Fine-tuning teaches \emph{style, syntax, and task behavior}; RAG supplies \emph{factual knowledge}.
    \item \textbf{Unchecked Direct Actuator Execution}: Connecting raw, unvalidated LLM text output directly to physical actuators (e.g., opening a high-pressure steam valve based on unconstrained text). Always pass LLM outputs through a deterministic software validation layer with hard safety interlocks before triggering physical hardware.
    \item \textbf{Using Excessive LoRA Rank ($r > 64$)}: Setting the LoRA rank unnecessarily high ($r = 128$ or $256$). A large rank increases parameters, slows down training, and causes severe overfitting on small domain datasets. Ranks of $r = 8$ or $r = 16$ achieve optimal performance for almost all edge adaptation tasks.
\end{enumerate}
\end{warning}

% ==============================================================================
% 9. KEY TAKEAWAYS
% ==============================================================================
\subsection{Key Takeaways}\label{subsec:takeaways}

\begin{takeaways}[Core Principles of Large Language Models]
\begin{itemize}
    \item \textbf{The 3-Stage Training Pipeline}: LLMs advance from self-supervised Pretraining (predicting next tokens) to Supervised Fine-Tuning (learning instruction following) to Alignment via RLHF/DPO (ensuring safety and helpfulness).
    \item \textbf{RAG Prevents Hallucinations}: Retrieval-Augmented Generation injects retrieved private documents into the prompt at query time, anchoring model outputs in factual truth.
    \item \textbf{LoRA Slashes Adaptation Costs}: Low-Rank Adaptation freezes the base model and trains tiny rank-decomposition matrices ($\Delta \mathbf{W} = \mathbf{B}\mathbf{A}$), slashing trainable parameters by over $99\%$.
    \item \textbf{Structured JSON for IoT}: Constrained grammar decoding forces LLMs to emit deterministic, valid JSON payloads that can directly control edge IoT actuators.
    \item \textbf{Small Language Models on the Edge}: 4-bit quantized SLMs (1B to 3B parameters) fit into under 2 GB of RAM, enabling local, private, and offline intelligence on industrial gateways.
\end{itemize}
\end{takeaways}

% ==============================================================================
% 10. CHECK YOURSELF
% ==============================================================================
\subsection{Check Yourself}\label{subsec:check_yourself}

\begin{checkyourself}[Practice Questions and Calculations]
\begin{enumerate}
    \item \textbf{LoRA Parameter Calculation}: A linear projection layer in a Transformer model has input dimension $k = 2{,}048$ and output dimension $d = 2{,}048$.
    \begin{enumerate}[label=(\alph*)]
        \item How many trainable weights are in this layer during full fine-tuning?
        \item If we apply a LoRA adapter with rank $r = 4$, how many trainable weights are in matrices $\mathbf{A}$ and $\mathbf{B}$ combined?
        \item What is the percentage parameter reduction?
    \end{enumerate}
    \item \textbf{Perplexity from Cross-Entropy Loss}: An evaluation run on an edge small language model yields an average cross-entropy loss of $\mathcal{L} = 2.3026$ nats per token. What is the model's perplexity ($\operatorname{PPL}$)?
    \item \textbf{Cosine Similarity Ranking in RAG}: A query vector is $\mathbf{q} = [1, 0]^T$. Two document chunk vectors are $\mathbf{d}_1 = [0.8, 0.6]^T$ and $\mathbf{d}_2 = [0.5, 0.5]^T$. Calculate the cosine similarity for both documents. Which document will the RAG system retrieve as the Top-1 match?
    \item \textbf{LoRA Initialization Property}: In LoRA, matrix $\mathbf{A}$ is initialized from a normal distribution $\mathcal{N}(0, \sigma^2)$, while matrix $\mathbf{B}$ is initialized to all zeros ($\mathbf{B} = \mathbf{0}$). What is the value of $\Delta \mathbf{W} = \mathbf{B}\mathbf{A}$ at the very beginning of training? Why is this mathematically vital?
    \item \textbf{Edge RAM Budget for 4-Bit SLMs}: An industrial edge gateway (such as a Raspberry Pi 5) has 4 GB of total system RAM, of which 1.5 GB is reserved for the operating system and video streams. Can a 3-billion parameter model quantized to 4-bit precision (requiring $\approx 1.8\text{ GB}$ of RAM) fit into the remaining available memory?
\end{enumerate}
\end{checkyourself}

% ==============================================================================
% 11. ANSWERS
% ==============================================================================
\subsection*{Answers}\label{subsec:answers}

\begin{enumerate}
    \item \textbf{LoRA Parameter Calculation}:
    \begin{enumerate}[label=(\alph*)]
        \item \textbf{Full Fine-Tuning}:
        \begin{equation*}
        \text{Weights}_{\text{full}} = d \times k = 2{,}048 \times 2{,}048 = \mathbf{4{,}194{,}304 \text{ parameters}}
        \end{equation*}

        \item \textbf{LoRA Adapter Parameters}:
        Matrix $\mathbf{A} \in \mathbb{R}^{r \times k}$ and Matrix $\mathbf{B} \in \mathbb{R}^{d \times r}$ with $r = 4$:
        \begin{align*}
        \text{Params}(\mathbf{A}) &= 4 \times 2{,}048 = 8{,}192 \\
        \text{Params}(\mathbf{B}) &= 2{,}048 \times 4 = 8{,}192 \\
        \text{Total LoRA Params} &= 8{,}192 + 8{,}192 = \mathbf{16{,}384 \text{ parameters}}
        \end{align*}

        \item \textbf{Percentage Reduction}:
        \begin{equation*}
        \text{Reduction} = \left(1 - \frac{16{,}384}{4{,}194{,}304}\right) \times 100\% = (1 - 0.003906) \times 100\% \approx \mathbf{99.61\%}
        \end{equation*}
        LoRA eliminates over $99.6\%$ of trainable parameters.
    \end{enumerate}

    \item \textbf{Perplexity Calculation}:
    Using Equation~\eqref{eq:perplexity}:
    \begin{equation*}
    \operatorname{PPL} = \exp(\mathcal{L}) = \exp(2.3026) = e^{2.3026} \approx \mathbf{10.0}
    \end{equation*}
    A perplexity of $10.0$ means the model is as uncertain about the next token as if it were choosing uniformly at random among 10 candidate words.

    \item \textbf{Cosine Similarity Ranking in RAG}:
    Using Equation~\eqref{eq:cosine_similarity}:
    \begin{itemize}
        \item For $\mathbf{d}_1 = [0.8, 0.6]^T$:
        \begin{align*}
        \|\mathbf{q}\|_2 &= \sqrt{1^2 + 0^2} = 1.0 \\
        \|\mathbf{d}_1\|_2 &= \sqrt{0.8^2 + 0.6^2} = \sqrt{0.64 + 0.36} = \sqrt{1.0} = 1.0 \\
        \mathbf{q} \cdot \mathbf{d}_1 &= (1 \times 0.8) + (0 \times 0.6) = 0.8 \\
        \cos(\theta_1) &= \frac{0.8}{1.0 \times 1.0} = \mathbf{0.8000}
        \end{align*}
        \item For $\mathbf{d}_2 = [0.5, 0.5]^T$:
        \begin{align*}
        \|\mathbf{d}_2\|_2 &= \sqrt{0.5^2 + 0.5^2} = \sqrt{0.25 + 0.25} = \sqrt{0.5} \approx 0.7071 \\
        \mathbf{q} \cdot \mathbf{d}_2 &= (1 \times 0.5) + (0 \times 0.5) = 0.5 \\
        \cos(\theta_2) &= \frac{0.5}{1.0 \times 0.7071} = \frac{0.5}{0.7071} \approx \mathbf{0.7071}
        \end{align*}
        \item \textbf{Conclusion}: $\cos(\theta_1) = 0.8000 > \cos(\theta_2) = 0.7071$. The RAG system selects \textbf{Document Chunk $\mathbf{d}_1$} as the Top-1 retrieved match.
    \end{itemize}

    \item \textbf{LoRA Initialization Property}:
    At step zero, since $\mathbf{B} = \mathbf{0}$:
    \begin{equation*}
    \Delta \mathbf{W} = \mathbf{B} \mathbf{A} = \mathbf{0} \times \mathbf{A} = \mathbf{0}
    \end{equation*}
    The total initial weight is $\mathbf{W} = \mathbf{W}_0 + \Delta \mathbf{W} = \mathbf{W}_0 + \mathbf{0} = \mathbf{W}_0$.
    This is mathematically vital because it ensures that before any fine-tuning gradient steps occur, \textbf{the model's behavior is 100\% identical to the original pretrained foundation model}. Training begins smoothly from the pretrained baseline without any initial performance drop or random noise shock.

    \item \textbf{Edge RAM Budget for 4-Bit SLMs}:
    \begin{align*}
    \text{Total Gateway RAM} &= 4.0 \text{ GB} \\
    \text{Reserved RAM (OS, Video, Buffers)} &= 1.5 \text{ GB} \\
    \text{Free Available RAM} &= 4.0 - 1.5 = \mathbf{2.5 \text{ GB}}
    \end{align*}
    Because the 4-bit 3B model requires $\approx 1.8\text{ GB}$ of RAM, and the available free memory is $2.5\text{ GB}$, the model will \textbf{comfortably fit into memory with approximately $700\text{ MB}$ of headroom} remaining for runtime context caches and active inference activations.
\end{enumerate}
